\input{preamble.tex}

\title{Лабораторная работа № 9. \\ Использование протокола STP. Агрегирование каналов}
\author{Сущенко Алина Николаевна \\ НПИбд-01-23}

\institute{Российский университет дружбы народов имени Патриса Лумумбы}
\date{2026}

\begin{document}

\frame{\titlepage}

\begin{frame}
\frametitle{Цель работы}
\begin{itemize}
    \item Изучение возможностей протокола STP и его модификаций по обеспечению отказоустойчивости сети.
    \item Агрегирование интерфейсов и перераспределение нагрузки между ними.
\end{itemize}
\end{frame}

\begin{frame}
\frametitle{Топология сети}
    \centering
    \includegraphics[width=0.7\textwidth]{../images/01.png}
    \captionof{figure}{Обновлённая топология сети.}
\end{frame}

\begin{frame}
\frametitle{Настройка порта Gig0/2 на msk-ansusenko-sw-01}
    \centering
    \includegraphics[width=\textwidth]{../images/02.png}
    \captionof{figure}{Настройка порта Gig0/2 на msk-ansusenko-sw-01.}
\end{frame}

\begin{frame}
\frametitle{Настройка порта Gig0/1 на msk-ansusenko-sw-4}
    \centering
    \includegraphics[width=\textwidth]{../images/03.png}
    \captionof{figure}{Настройка порта Gig0/1 на msk-ansusenko-sw-4.}
\end{frame}

\begin{frame}
\frametitle{Настройка транкового режима}
    \centering
    \includegraphics[width=\textwidth]{../images/04.png}
    \captionof{figure}{Настройка транкового режима.}
\end{frame}

\begin{frame}
\frametitle{Настройка транкового режима на порту Fa0/23}
    \centering
    \includegraphics[width=\textwidth]{../images/05.png}
    \captionof{figure}{Настройка транкового режима на порту Fa0/23.}
\end{frame}

\begin{frame}
\frametitle{Настройка транков на msk-ansusenko-sw-4}
    \centering
    \includegraphics[width=\textwidth]{../images/06.png}
    \captionof{figure}{Настройка транков на msk-ansusenko-sw-4.}
\end{frame}

\begin{frame}
\frametitle{Проверка статуса транков на msk-ansusenko-sw-01}
    \centering
    \includegraphics[width=\textwidth]{../images/07.png}
    \captionof{figure}{Проверка статуса транков на msk-ansusenko-sw-01.}
\end{frame}

\begin{frame}
\frametitle{Проверка статуса транков на msk-ansusenko-sw-3}
    \centering
    \includegraphics[width=\textwidth]{../images/08.png}
    \captionof{figure}{Проверка статуса транков на msk-ansusenko-sw-3.}
\end{frame}

\begin{frame}
\frametitle{Проверка статуса транков на msk-ansusenko-sw-4}
    \centering
    \includegraphics[width=\textwidth]{../images/09.png}
    \captionof{figure}{Проверка статуса транков на msk-ansusenko-sw-4.}
\end{frame}

\begin{frame}
\frametitle{Пинг серверов с dk-ansusenko-01}
    \centering
    \includegraphics[width=\textwidth]{../images/10.png}
    \captionof{figure}{Пинг серверов с dk-ansusenko-01.}
\end{frame}

\begin{frame}
\frametitle{Пинг почтового сервера через коммутатор}
    \centering
    \includegraphics[width=\textwidth]{../images/11.png}
    \captionof{figure}{Пинг почтового сервера через коммутатор msk-ansusenko-sw-2.}
\end{frame}

\begin{frame}
\frametitle{Состояние STP для VLAN 3 на msk-ansusenko-sw-2}
    \centering
    \includegraphics[width=\textwidth]{../images/12.png}
    \captionof{figure}{Состояние STP для VLAN 3 на msk-ansusenko-sw-2.}
\end{frame}

\begin{frame}
\frametitle{Настройка корневого коммутатора}
    \centering
    \includegraphics[width=\textwidth]{../images/13.png}
    \captionof{figure}{Настройка корневого коммутатора.}
\end{frame}

\begin{frame}
\frametitle{Состояние для VLAN 3 на корневом коммутаторе}
    \centering
    \includegraphics[width=\textwidth]{../images/14.png}
    \captionof{figure}{Состояние для VLAN 3 на корневом коммутаторе.}
\end{frame}

\begin{frame}
\frametitle{Движение пакетов до mail}
    \centering
    \includegraphics[width=\textwidth]{../images/15.png}
    \captionof{figure}{Движение пакетов ICMP от dk-ansusenko-01 до mail через msk-ansusenko-sw-01 и msk-ansusenko-sw-3.}
\end{frame}

\begin{frame}
\frametitle{Движение пакетов до web}
    \centering
    \includegraphics[width=\textwidth]{../images/16.png}
    \captionof{figure}{Движение пакетов ICMP от dk-ansusenko-01 до web через msk-ansusenko-sw-01 и msk-ansusenko-sw-2.}
\end{frame}

\begin{frame}
\frametitle{Настройка Portfast на msk-ansusenko-sw-2}
    \centering
    \includegraphics[width=\textwidth]{../images/17.png}
    \captionof{figure}{Настройка Portfast на msk-ansusenko-sw-2.}
\end{frame}

\begin{frame}
\frametitle{Настройка Portfast на msk-ansusenko-sw-3}
    \centering
    \includegraphics[width=\textwidth]{../images/18.png}
    \captionof{figure}{Настройка Portfast на msk-ansusenko-sw-3.}
\end{frame}

\begin{frame}
\frametitle{Длительный пинг до почтового сервера (с разрывом)}
    \centering
    \includegraphics[width=\textwidth]{../images/19.png}
    \captionof{figure}{Длительный пинг до почтового сервера с разрывом соединения.}
\end{frame}

\begin{frame}
\frametitle{Длительный пинг до почтового сервера (без разрыва)}
    \centering
    \includegraphics[width=\textwidth]{../images/20.png}
    \captionof{figure}{Длительный пинг до почтового сервера без разрыва соединения.}
\end{frame}

\begin{frame}
\frametitle{Логическая схема с агрегированным соединением}
    \centering
    \includegraphics[width=\textwidth]{../images/21.png}
    \captionof{figure}{Логическая схема локальной сети с агрегированным соединением.}
\end{frame}

\begin{frame}
\frametitle{Настройка EtherChannel на msk-ansusenko-sw-01}
    \centering
    \includegraphics[width=\textwidth]{../images/22.png}
    \captionof{figure}{Настройка EtherChannel для msk-ansusenko-sw-01.}
\end{frame}

\begin{frame}
\frametitle{Настройка EtherChannel на msk-ansusenko-sw-4}
    \centering
    \includegraphics[width=\textwidth]{../images/23.png}
    \captionof{figure}{Настройка EtherChannel для msk-ansusenko-sw-4.}
\end{frame}

\begin{frame}
\frametitle{Контрольный вопрос №1}
\begin{block}{Какую информацию можно получить, воспользовавшись командой определения состояния протокола STP для VLAN (на корневом и не на корневом устройстве)? Приведите примеры.}
Команда \texttt{show spanning-tree vlan <id>} позволяет получить:
\begin{itemize}
\item Приоритет и MAC-адрес корневого устройства (Root ID)
\item Приоритет и MAC-адрес текущего моста (Bridge ID)
\item Роли портов (Root, Designated, Alternate)
\item Состояние портов (FWD, BLK, LSN, LRN)
\item Стоимость портов (Cost) и приоритеты портов (Prio.Nbr)
\end{itemize}
\end{block}
\end{frame}

\begin{frame}[containsverbatim]
\frametitle{Контрольный вопрос №1 (пример вывода)}
\begin{block}{На корневом устройстве (msk-ansusenko-sw-01):}
\begin{minted}{text}
Root ID    Priority    24579
           Address     00D0.BCEE.A8B6
           This bridge is the root
Bridge ID  Priority    24579
           Address     00D0.BCEE.A8B6
Interface        Role Sts Cost
Fa0/1            Desg FWD 19
Gi0/1            Desg FWD 4
Gi0/2            Desg FWD 1
\end{minted}
\end{block}
\end{frame}

\begin{frame}[containsverbatim]
\frametitle{Контрольный вопрос №2}
\begin{block}{При помощи какой команды можно узнать, в каком режиме (STP или Rapid PVST+) работает устройство?}
Команда \texttt{show spanning-tree vlan <id>}. В первой строке вывода указывается используемый протокол:
\begin{itemize}
\item \texttt{Spanning tree enabled protocol ieee} - классический STP (802.1d)
\item \texttt{Spanning tree enabled protocol rstp} - Rapid STP (802.1w)
\item \texttt{Spanning tree enabled protocol mstp} - Multiple STP (802.1s)
\end{itemize}
\end{block}
\end{frame}

\begin{frame}[containsverbatim]
\frametitle{Контрольный вопрос №2 (пример вывода)}
\begin{block}{Пример вывода для STP:}
\begin{minted}{text}
VLAN0003
  Spanning tree enabled protocol ieee
  Root ID    Priority    32771
             Address     0001.4272.C298
             This bridge is the root
\end{minted}
\end{block}
\begin{block}{Пример вывода для Rapid PVST+:}
\begin{minted}{text}
VLAN0003
  Spanning tree enabled protocol rstp
  Root ID    Priority    24579
             Address     00D0.BCEE.A8B6
\end{minted}
\end{block}
\end{frame}

\begin{frame}
\frametitle{Контрольный вопрос №3}
\begin{block}{Для чего и в каких случаях нужно настраивать режим Portfast?}
Portfast настраивается на access-портах коммутатора, к которым подключены конечные устройства (компьютеры, серверы, принтеры).

\textbf Позволяет порту сразу переходить в режим Forwarding, минуя стадии Listening и Learning (15 + 15 = 30 секунд).

\textbf{Когда использовать:}
\begin{itemize}
\item На портах, подключенных к одному хосту
\item На портах, где не может образоваться петля
\item Для ускорения загрузки ОС по сети (PXE)
\end{itemize}

\end{block}
\end{frame}

\begin{frame}
\frametitle{Контрольный вопрос №4}
\begin{block}{В чем состоит принцип работы агрегированного интерфейса? Для чего он используется?}
\textbf EtherChannel объединяет несколько физических портов в один логический канал (Port-channel).

\textbf{Особенности:}
\begin{itemize}
\item Трафик распределяется между физическими портами по алгоритму хэширования (MAC-адреса, IP-адреса, порты TCP/UDP)
\item Все порты в группе должны иметь одинаковые настройки (скорость, дуплекс, VLAN, trunk)
\end{itemize}

\textbf{Для чего используется:}
\begin{itemize}
\item Увеличение пропускной способности (агрегирование полосы)
\item Отказоустойчивость (при выходе порта из строя трафик идёт по остальным)
\item Балансировка нагрузки между портами
\end{itemize}
\end{block}
\end{frame}

\begin{frame}
\frametitle{Контрольный вопрос №5}
\begin{block}{В чём принципиальные отличия при использовании протоколов LACP, PAgP и статического агрегирования?}
\begin{tabular}{|p{2.5cm}|p{3.5cm}|p{3cm}|p{2.5cm}|}
\hline
\textbf{Характеристика} & \textbf{LACP (802.3ad)} & \textbf{PAgP (Cisco)} & \textbf{Static (mode on)} \\
\hline
Тип & Открытый стандарт & Проприетарный Cisco & Без протокола \\
\hline
Совместимость & Разные вендоры & Только Cisco & Любое оборудование \\
\hline
Обнаружение ошибок & Да & Да & Нет \\
\hline
Режимы & active/passive & desirable/auto & on \\
\hline
\end{tabular}

\textbf{Вывод:} LACP рекомендуется для мультивендорных сетей, PAgP - для сетей только Cisco, Static - для простых конфигураций без необходимости протокольного согласования.
\end{block}
\end{frame}

\begin{frame}[containsverbatim]
\frametitle{Контрольный вопрос №6}
\begin{block}{При помощи каких команд можно узнать состояние агрегированного канала EtherChannel?}
Основные команды для проверки EtherChannel:
\begin{minted}{bash}
# Краткая информация о всех EtherChannel
show etherchannel summary

# Детальная информация о Port-channel
show etherchannel port-channel

# Статус логического интерфейса
show interfaces port-channel <номер>

# Статус физических портов в группе
show interfaces status

# Информация о балансировке нагрузки
show etherchannel load-balance
\end{minted}
\end{block}
\end{frame}

\begin{frame}[containsverbatim]
\frametitle{Контрольный вопрос №6 (пример вывода)}
\begin{block}{Пример вывода \texttt{show etherchannel summary}:}
\begin{minted}{text}
Flags:  D - down        P - in port-channel
        I - stand-alone s - suspended
Group  Port-channel  Protocol  Ports
------+-------------+---------+----------------
1      Po1(SU)        -       Fa0/20(P) Fa0/21(P)
                              Fa0/22(P) Fa0/23(P)
\end{minted}
\textbf{Расшифровка:} SU - Port-channel работает (S - switched, U - up), все порты в группе активны (P).
\end{block}
\end{frame}

\begin{frame}
\frametitle{Выводы}
\begin{itemize}
	\item В ходе выполнения лабораторной работы были изучены протоколы STP и Rapid PVST+, а также технология агрегирования каналов EtherChannel. Настроено резервное соединение между коммутаторами, обеспечена отказоустойчивость сети, активирован режим Portfast для серверных портов. Проверена работоспособность ICMP-запросов в режиме симуляции. Все поставленные задачи выполнены.
\end{itemize}
\end{frame}

\end{document}
